\documentclass{article}
\usepackage[T1]{fontenc}
\usepackage{lmodern}
\usepackage{graphicx}
\usepackage{amsmath,amssymb}
\usepackage[a4paper,margin=2cm]{geometry}
\usepackage[absolute,overlay]{textpos}
\setlength{\TPHorizModule}{1mm}
\setlength{\TPVertModule}{1mm}
\usepackage[indonesian]{babel}
\usepackage{hyperref}
\setlength{\emergencystretch}{2em}
% unit: Halaman 56

\begin{document}

% === Halaman 56 (llm) ===

\section*{2. Diffusion}

\begin{itemize}
    \item Prinsip ini menyebarkan pengaruh satu bit plainteks atau kunci ke sebanyak mungkin bit cipherteks.
    Contoh: satu huruf plainteks mempengaruhi nilai banyak huruf cipherteks
    
    \item Sebagai efek difusi, maka pengubahan kecil pada plainteks sebanyak satu atau dua bit menghasilkan perubahan pada bit-bit cipherteks yang tidak dapat diprediksi.
    
    \item Contoh: mengenkripsi plainteks \texttt{1111111111111111} \\ 
    menghasilkan cipherteks \texttt{01101100000101001}. \\ 
    Misalkan 1 bit pertama plainteks diubah menjadi \texttt{0111111111111111}. Jika difusi bekerja, maka cipherteksnya berubah signifikan menjadi \texttt{10110010111011111}
    
    \item \textit{Diffusion} did alam block cipher dapat direalisasikan dengan menggunakan teknik permutasi atau transposisi secara berulang-ulang. Permutasi adalah mengacak susunan bit atau \textit{byte} di dalam plainteks.
\end{itemize}

\end{document}
